\vspace{-1em}
\section{Experiments}
\vspace{-0.5em}
\subsection{Experimental Settings}

\textbf{Datasets.} We trained OmniSync using the MEAD dataset~\cite{wang2020mead} and a 400-hour dataset collected from YouTube. MEAD's controlled laboratory recordings with diverse facial expressions but minimal head movement provided ideal data for training early denoising stages, while the YouTube dataset enhanced generalization across varied real-world conditions for middle and late stages.

To address the limitations of existing benchmarks that focus solely on real-world videos with frontal views and stable lighting, we created the AIGC-LipSync Benchmark. This comprehensive evaluation framework comprises 615 human-centric videos generated by state-of-the-art text-to-video models such as Kling, Dreamina, Wan~\cite{wang2025wan}, and Hunyuan~\cite{kong2024hunyuanvideo}. The benchmark specifically captures challenging visual scenarios such as large facial movements, profile views, variable lighting, occlusions, and stylized characters—conditions that traditional benchmarks fail to address. Details about benchmark construction can be found in the supplementary materials.

\textbf{Implementation Details.} We implement our OmniSync model using the Diffusion Transformer architecture. The model is trained on a combined dataset for 80,000 steps using AdamW optimizer~\cite{loshchilov2017decoupled} with a learning rate of 1e-5. Training is completed in 80 hours using 64 NVIDIA A100 GPUs with a batch size of 64. Audio features are extracted via a pre-trained Whisper encoder, and text conditioning utilizes a T5 encoder. Training employs the timestep‑dependent sampling threshold $t_{\text{threshold}}= 850$. The experimental results indicate that excessive thresholds induce significant misalignment while insufficient values will leak the original lip shape.
During inference we adopt our flow‑matching‑based progressive noise initialization starting at $\tau = 0.92$,  followed by 50 denoising steps.




\begin{table}[]
\caption{Quantitative comparison with previous methods on HDTF Dataset.}
\label{tab:results1}
\resizebox{\linewidth}{!}{
\setlength{\tabcolsep}{4pt}
\begin{tabular}{@{}l ccc ccc cc@{}}
\toprule
\multicolumn{9}{c}{\textbf{HDTF Dataset}} \\
\midrule
\multirow{2}{*}[-0.5ex]{\textbf{Method}} & \multicolumn{3}{c}{\textbf{Full Reference Metrics}} & \multicolumn{3}{c}{\textbf{No Reference Metrics}} & \multicolumn{2}{c}{\textbf{Lip Sync}} \\
\cmidrule(lr){2-4} \cmidrule(lr){5-7} \cmidrule(lr){8-9}
& FID  & FVD  & CSIM $\uparrow$ & NIQE  & BRISQUE  & HyperIQA $\uparrow$ & LMD  & LSE-C $\uparrow$ \\
\midrule
Wav2Lip~\cite{prajwal2020lip}        & 14.912 & 543.340 & 0.852 & 6.495 & 53.372 & 45.822 & 10.007 & 7.630 \\
VideoReTalking~\cite{cheng2022videoretalking} & 11.868 & 379.518 & 0.786 & 6.333 & 50.722 & 48.476 & 8.848  & 7.180 \\
TalkLip~\cite{wang2023seeing}        & 16.680 & 691.518 & 0.843 & 6.377 & 52.109 & 44.393 & 15.954 & 5.880 \\
IP-LAP~\cite{zhong2023identity}         & 9.512  & 325.691 & 0.809 & 6.533 & 54.402 & 50.086 & 7.695  & 7.260 \\
Diff2Lip~\cite{mukhopadhyay2024diff2lip}       & 12.079 & 461.341 & 0.869 & 6.261 & 49.361 & 48.869 & 18.986 & 7.140 \\
MuseTalk~\cite{zhang2024musetalk}       & 8.759  & 231.418 & 0.862 & 5.824 & 46.003 & 55.397 & 8.701  & 6.890 \\
LatentSync~\cite{li2024latentsync}     & 8.518  & 216.899 & 0.859 & 6.270 & 50.861 & 53.208 & 17.344 & \textbf{8.050} \\
\textbf{Ours}  & \textbf{7.855} & \textbf{199.627} & \textbf{0.875} & \textbf{5.481} & \textbf{37.917} & \textbf{56.356} & \textbf{7.097} & 7.309 \\
\bottomrule
\end{tabular}}
\vspace{-0.5em}
\end{table}


\begin{table}[]
\caption{Quantitative comparison with previous methods on AIGC-LipSync Benchmark.}
% \vspace{-1em}
\label{tab:results2}
\resizebox{\linewidth}{!}{
% \centering
% 
\setlength{\tabcolsep}{4pt}
\begin{tabular}{@{}l ccc ccc cc@{}}
\toprule
\multicolumn{9}{c}{\textbf{AIGC-LipSync Benchmark}} \\
\midrule
\multirow{2}{*}[-1.5ex]{\textbf{Method}} & \multicolumn{3}{c}{\textbf{Full Reference Metrics}} & \multicolumn{3}{c}{\textbf{No Reference Metrics}} & \multicolumn{2}{c}{\textbf{Generation Success Rate}} \\
\cmidrule(lr){2-4} \cmidrule(lr){5-7} \cmidrule(lr){8-9}
& FID  & FVD  & CSIM $\uparrow$ & NIQE  & BRISQUE  & HyperIQA $\uparrow$ & \begin{tabular}[c]{@{}c@{}}All Videos $\uparrow$\end{tabular} & \begin{tabular}[c]{@{}c@{}}Stylized\\Characters $\uparrow$\end{tabular} \\
\midrule
Wav2Lip~\cite{prajwal2020lip}        & 22.989 & 562.245 & 0.727 & 5.392 & 42.816 & 50.511 & 71.38\% & 26.67\% \\
VideoReTalking~\cite{cheng2022videoretalking} & 20.439 & 329.460 & 0.669 & 5.947 & 45.047 & 48.645 & 48.78\% & 7.78\% \\
TalkLip~\cite{wang2023seeing}        & 31.180 & 619.179 & 0.754 & 5.239 & 41.692 & 50.608 & 52.36\% & 34.44\% \\
IP-LAP~\cite{zhong2023identity}         & 14.686 & 247.402 & 0.796 & 5.546 & 45.153 & 53.174 & 45.53\% & 6.67\% \\
Diff2Lip~\cite{mukhopadhyay2024diff2lip}       & 23.542 & 403.149 & 0.692 & 5.440 & 42.442 & 50.335 & 74.63\% & 36.67\% \\
MuseTalk~\cite{zhang2024musetalk}       & 17.668 & 297.621 & 0.667 & 4.935 & 36.017 & 58.334 & 92.20\% & 67.78\% \\
LatentSync~\cite{li2024latentsync}     & 15.374 & 263.111 & 0.751 & 5.342 & 41.917 & 54.648 & 74.96\% & 35.56\% \\
\textbf{Ours}  & \textbf{10.681} & \textbf{211.350} & \textbf{0.808} & \textbf{4.588} & \textbf{25.485} & \textbf{61.906} & \textbf{97.40\%} & \textbf{87.78\%} \\
\bottomrule
\end{tabular}}
\vspace{-1em}
\end{table}



\vspace{-0.5em}
\subsection{Quantitative Evaluation}
\vspace{-0.5em}

We evaluate OmniSync against state-of-the-art methods including Wav2Lip~\cite{prajwal2020lip}, VideoReTalking~\cite{cheng2022videoretalking}, TalkLip~\cite{wang2023seeing}, IP-LAP~\cite{zhong2023identity}, Diff2Lip~\cite{mukhopadhyay2024diff2lip}, MuseTalk~\cite{zhang2024musetalk}, and LatentSync~\cite{li2024latentsync} using a comprehensive suite of metrics. For visual quality assessment, we employ FID (Fréchet Inception Distance) to measure frame-level fidelity, FVD (Fréchet Video Distance) for temporal consistency, and CSIM (Cosine Similarity) to quantify identity preservation. Perceptual quality is assessed using no-reference metrics including NIQE (Natural Image Quality Evaluator), BRISQUE (Blind/Referenceless Image Spatial Quality Evaluator), and HyperIQA~\cite{su2020blindly}. For audio-visual synchronization, we measure LMD (Landmark Distance) between predicted and ground truth facial landmarks in the mouth region, and LSE-C (Lip Sync Error - Confidence) to evaluate lip synchronization quality.

For the AIGC-LipSync benchmark, we report the Generation Success Rate across all 615 videos and specifically for stylized characters. This metric shows the percentage of videos that are successfully synchronized and pass human verification. This evaluation is essential for universal lip synchronization in AI-generated content, where traditional metrics may not fully capture the challenges of stylized characters, extreme poses, and other atypical visual conditions.


The experimental results in Tab.~\ref{tab:results1} and Tab.~\ref{tab:results2} demonstrate that our approach achieves superior performance on multiple metrics. On the HDTF dataset, our method reduced FID by 7.8$\%$ and FVD by 8.0$\%$ compared to LatentSync, with a remarkable 23.2$\%$ improvement in BRISQUE over Diff2Lip. For lip synchronization, we achieved the lowest LMD, outperforming IP-LAP by 7.8$\%$, while LatentSync maintained a slight edge in LSE-C due to its SyncNet-based loss constraint.

On the challenging AIGC-LipSync benchmark, OmniSync demonstrated exceptional capabilities with a 30.5$\%$ FID reduction and 19.7$\%$ FVD reduction compared to LatentSync, alongside improved identity preservation. Most significantly, our method achieved a 97.40$\%$ Generation Success Rate across all videos—substantially higher than MuseTalk (92.20$\%$) and other methods (below 75$\%$). For stylized characters, our success rate of 87.78$\%$ outperformed MuseTalk (67.78$\%$), demonstrating OmniSync's capability to handle diverse visual representations including stylized characters.

\begin{figure*}
\vspace{-1em}
\begin{center}
   \includegraphics[width=1.\linewidth]{imgs/omnisync-comparsion-0516_compressed.pdf}
\end{center}
\vspace{-1em}
   \caption{\textbf{Qualitative comparison} with previous methods across diverse subjects and phonemes. Our approach produces more accurate lip synchronization and better identity preservation.}
\label{fig:com}
\vspace{-1em}
\end{figure*}

\begin{table}[]
\caption{\textbf{User study} results comparing various audio-driven lip-sync methods.}
\label{tab:user_study}
\resizebox{\linewidth}{!}{
\begin{tabular}{@{}lcccccccc@{}}
\toprule
\textbf{Metric}    & \textbf{Wav2Lip} & \textbf{\begin{tabular}[c]{@{}c@{}}Video\\ ReTalking\end{tabular}} & \textbf{TalkLip} & \textbf{IP-LAP} & \textbf{Diff2Lip} & \textbf{MuseTalk} & \textbf{\begin{tabular}[c]{@{}c@{}}Latent\\ Sync\end{tabular}} & \textbf{Ours}  \\ \midrule
Lip Sync Accuracy  & 2.684            & 2.769                                                              & 1.940            & 2.359           & 3.000             & 2.632             & 3.812                                                          & \textbf{3.923} \\
Identity Preservation & 2.410            & 2.786                                                              & 2.222            & 2.991           & 2.889             & 3.197             & 3.658                                                          & \textbf{4.128} \\
Timing Stability   & 2.556            & 2.376                                                              & 2.162            & 3.034           & 2.812             & 3.145             & 3.581                                                          & \textbf{4.043} \\
Image Quality      & 2.017            & 2.607                                                              & 1.889            & 3.171           & 2.402             & 3.094             & 3.632                                                          & \textbf{4.051} \\
Video Realism      & 2.120            & 2.419                                                              & 1.838            & 2.316           & 2.479             & 2.761             & 3.453                                                          & \textbf{3.872} \\ \bottomrule
\end{tabular}}
% \vspace{-1em}
\end{table}



\subsection{Qualitative Evaluation}
We present qualitative comparisons between OmniSync and existing methods in Fig.~\ref{fig:com}. Our approach produces more natural facial expressions and superior lip synchronization. Due to lip shape leakage, IP-LAP~\cite{zhong2023identity} and LatentSync~\cite{li2024latentsync} frequently fail at mouth shape modification, resulting in poor lip synchronization effects. MuseTalk~\cite{zhang2024musetalk} and VideoReTalking~\cite{cheng2022videoretalking} modify lip movements but frequently lose identity and visual quality. Diff2Lip~\cite{mukhopadhyay2024diff2lip} and Wav2Lip~\cite{prajwal2020lip} commonly exhibit lip sync errors, mouth artifacts, and identity drift, particularly in challenging or stylized cases. In contrast, OmniSync consistently maintains identity details and generates realistic, expressive lip movements, demonstrating robust performance. Our approach effectively balances audio and visual cues, addressing the challenge of weak audio conditioning.

\textbf{User Study.}
To assess perceptual quality, we conducted a user study with 39 participants evaluating 32 video sets generated by OmniSync and seven competing methods, with a standardized Cronbach's $\alpha$ coefficient of 0.98. Participants rated each video on a 5-point Likert scale across five criteria: Lip Sync Accuracy, Character Identity preservation, Timing Stability, Image Quality, and Video Realism. 
As shown in Tab.~\ref{tab:user_study}, OmniSync outperformed all competitors across all metrics, achieving superior scores in Lip Sync Accuracy (3.923 vs. 3.812 for LatentSync), Character Identity (4.128 vs. 3.658), Timing Stability (4.043 vs. 3.581), Image Quality (4.051 vs. 3.632), and Video Realism (3.872 vs. 3.453). These results confirm OmniSync's superior ability to generate high-quality talking videos.



\subsection{Ablation Study}

To clarify the contributions of each core component in our framework, we conduct an ablation study targeting three key modules: the timestep-dependent data sampling strategy, progressive noise initialization, and the Dynamic Spatiotemporal Classifier-Free Guidance (DS-CFG) mechanism. Quantitative results are presented in Tab.~\ref{fig:ablation}, and corresponding visual examples are shown in Fig.~\ref{fig:ablation}. 

Removing the timestep-dependent sampling strategy results in a significant drop in identity preservation and pose consistency, with a 10.7$\%$ decrease in CSIM and substantial increases in FID and FVD. As shown in Fig.~\ref{fig:ablation}, without this sampling strategy, the generated faces often exhibit clear mismatches with the original image, including noticeable facial misalignment issues. This validates our design choice of aligning pseudo-paired data with early diffusion steps, which proves critical for generating structurally stable outputs. Similarly, removing progressive noise initialization leads to evident temporal inconsistencies and an increase in FVD, confirming the importance of our flow-matching initialization in preserving spatial anchoring and motion coherence. 

We also compare our proposed DS-CFG with both low and high static CFG settings. As illustrated in Fig.~\ref{fig:ablation}, low CFG provides insufficient audio conditioning, resulting in under-articulated lip movements (LSE-C: 4.16), whereas high CFG improves synchronization (LSE-C: 7.10) but introduces noticeable artifacts and distortions in facial details. In contrast, DS-CFG achieves an optimal balance by applying strong localized guidance in early diffusion stages and gradually reducing it in later steps. These results confirm that dynamic control across temporal and spatial dimensions is essential for producing expressive and visually coherent lip synchronization in generative video content.
% Please add the following required packages to your document preamble:
% \usepackage{booktabs}
% \usepackage{graphicx}
\begin{table}[]
\caption{\textbf{Ablation study} for our method.}
\label{tab:ablation}
\resizebox{\columnwidth}{!}{%
\begin{tabular}{@{}lccccccc@{}}
\toprule
Methods                                                                         & FID ↓           & FVD ↓            & CSIM ↑         & NIQE ↓         & BRISQUE ↓       & HyperIQA ↑      & LSE-C ↑       \\ \midrule
Ours                                                                            & \textbf{15.710} & \textbf{287.168} & \textbf{0.814} & \textbf{5.321} & \textbf{29.588} & \textbf{57.288} & 7.06          \\
\begin{tabular}[c]{@{}l@{}}w/o Timestep-Dependent\\ Sampling Strategy\end{tabular} & 21.552          & 549.768          & 0.727          & 5.462          & 30.346          & 56.204          & 7.00          \\
\begin{tabular}[c]{@{}l@{}}w/o Progressive Noise \\ Initialization\end{tabular} & 16.731          & 361.282          & 0.805          & 5.349          & 29.789          & 56.511          & 7.03          \\
w/ Low Static CFG                                                                      & -               & -                & -              & 5.359          & 29.724          & 56.568          & 4.16          \\
w/ High Static CFG                                                                     & 22.725          & 348.335          & 0.782          & 5.473          & 29.678          & 56.289          & \textbf{7.10} \\ \bottomrule
\end{tabular}}
\end{table}




\begin{figure*}
\vspace{-1em}
\begin{center}
   \includegraphics[width=1.\linewidth]{imgs/omnisync-ablation-0516_compressed.pdf}
\end{center}
\vspace{-1em}
   \caption{\textbf{Ablation study} for timestep-dependent sampling strategy and different CFG settings.}
\label{fig:ablation}
\vspace{-0.5em}
\end{figure*}

